\chapter{Discusión y análisis de resultados}
\label{cap:resultados}

% =====================================================================================
% GUION. Escrito en sep-2026 sobre el conjunto de prueba ya evaluado.
%
% 🔴 ORIGEN DE CADA CIFRA. Todas salen de estos tres ficheros y de ningun otro:
%      §8.1  -> notebooks/05_comparativa/05c_comparativa_escala.ipynb
%      §8.2  -> notebooks/05_comparativa/05d_mitigacion.ipynb
%      reglas -> notebooks/05_comparativa/METRICAS_Y_EVALUACION.md
%
% ⚠️ EL GUION ANTERIOR DECIA que toda interpretacion se fuera «al capitulo 11». Ese
%    capitulo no existe: la memoria tiene nueve. Y el titulo del capitulo —y el de la
%    estructura oficial del tipo 3— es «Resultados y discusion». Asi que la discusion va
%    aqui, pegada a la cifra que la sostiene, y las conclusiones del capitulo 9 se quedan
%    con lo que responde a los objetivos.
%
% ⚠️ FIGURAS. Las genera 05c/05d y se copian solas a figuras/cap08_v3/ y
%    figuras/cap08_mitigacion/. Ya vienen SIN titulo interno: lo que llevan arriba son
%    rotulos de panel. No regenerarlas a mano.
%
% ⚠️ De las 18 figuras y 15 tablas que exporta 05c entran CUATRO figuras y dos tablas
%    compactas, por decision de extension. El resto —MAE y Spearman de los tres ejes, las
%    vistas en diferencias y las tablas por dia— se queda en el cuaderno; si hiciera falta
%    reponer algo, el sitio es un anexo, no este capitulo.
% =====================================================================================

Este capítulo mide lo que los tres modelos del capítulo~\ref{cap:comparativa} hacen sobre el
conjunto de prueba. Se separa en dos partes: primero lo que el
modelo \textbf{aprende} (el factor de supervivencia, predicho antes de ejecutar) y después lo
que el modelo \textbf{sirve}, es decir, cuánto se acerca al valor exacto el observable una vez
corregido.

Además, se reportan el porcentaje de circuitos que se ha salvado de la \textbf{cota inferior del factor de supervivencia}.


% =====================================================================================
\section{Resultados sobre el factor de supervivencia}
\label{sec:res_factor}
% =====================================================================================

Antes de mirar la prueba conviene revisar la medida en validación
Tabla~\ref{tab:escalera_val}: mismos tipos de circuito y mismo rango de tamaños que el entrenamiento.

% ⚠️ FORMATO: titulo arriba y fuente abajo (Instrucciones §1.3, p. 6).
% Cifras: 05c, celda de `cmp.tabla_escalera`, columna «R² val». Redondeadas a dos
% decimales a proposito: la tercera no separa a ningun par de modelos.
\begin{table}[htbp]
    \centering
    \caption{Coeficiente de determinación sobre el factor, en validación.}
    \label{tab:escalera_val}
    {\footnotesize
    Fuente: elaboración propia. \begin{tabular}{@{}lrrrrr@{}}
        \toprule
        & \textbf{media $Z$} & \textbf{media $X$} & \textbf{media $Y$}
        & \textbf{paridad} & \textbf{corr. vec.} \\
        \midrule
        No mitigar            & $-0{,}43$ & $-0{,}61$ & $-0{,}53$ & $-1{,}78$ & $-0{,}84$ \\
        Factor constante      & $-0{,}00$ & $-0{,}00$ & $-0{,}00$ & $-0{,}01$ & $-0{,}01$ \\
        Ridge · 1 variable    & $0{,}39$ & $0{,}49$ & $0{,}44$ & $0{,}62$ & $0{,}59$ \\
        Ridge · 34 variables  & $0{,}41$ & $0{,}52$ & $0{,}45$ & $0{,}65$ & $0{,}60$ \\
        \textbf{Random Forest} & $\bm{0{,}42}$ & $\bm{0{,}53}$ & $\bm{0{,}46}$ & $\bm{0{,}67}$ & $\bm{0{,}61}$ \\
        GNN · SAGE            & $0{,}41$ & $0{,}52$ & $0{,}46$ & $0{,}64$ & $0{,}60$ \\
        GNN · GIN             & $0{,}41$ & $0{,}51$ & $0{,}45$ & $0{,}63$ & $0{,}59$ \\
        GNN · GATv2           & $0{,}40$ & $0{,}52$ & $0{,}45$ & $0{,}60$ & $0{,}57$ \\
        \bottomrule
    \end{tabular}}

    {\footnotesize Entre $2{.}282$ y $2{.}519$ circuitos por
    celda, según el observable.}
\end{table}

La escalera dice tres cosas. La primera, que \textbf{no mitigar es peor que
no hacer nada}, de modo que el listón de verdad no es el cero sino un factor constante. La segunda,
que \textbf{casi todo lo aprendible se aprende con una sola variable}: la duración total del
circuito ya consigue entre el $39$\,\% y el $62$\,\% de la métrica, y las treinta y tres
restantes añaden tres puntos como mucho (y en el correlador de vecinos, tres milésimas). Y la tercera, que \textbf{dentro de la distribución
de entrenamiento los cinco modelos empatan} y aunque Random Forest gana por un margen que no se ve en una gráfica.

Para desempatar entre los 5 modelos se comprobará en el conjunto de prueba.

\subsection{Eje temporal: la deriva del hardware}
\label{sec:res_tiempo}

La deriva se mide \textbf{en limpio}: solo circuitos aleatorios y de hasta once qubits, de modo
que lo único que varía entre celdas es el día de calibración
(Figura~\ref{fig:res_deriva}). Sin fijar el tipo y el tamaño, lo que se estaría midiendo es
otro eje.

% Figura ya exportada por 05c (v3_deriva_r2). Es la medida LIMPIA de deriva; la version
% mezclada (v3_tiempo_r2) trae ademas tipos y tamaños nuevos y no es interpretable.
\begin{figure}[htbp]
    \centering
    \caption{Coeficiente de determinación sobre el factor, por día de calibración, sobre
             circuitos aleatorios de hasta once qubits.}
    \label{fig:res_deriva}

    \includegraphics[width=\textwidth]{cap08_v3/v3_deriva_r2}

    {\footnotesize Fuente: elaboración propia. Entre $29$ y $60$ circuitos por celda.}
\end{figure}

El resultado es que \textbf{no hay deriva medible en doce días}: los cinco modelos suben y bajan juntos, con diferencias
entre ellos.

El modelo que destaca un poco por encima de un observable en concreto (paridad) es Random Forest, pero no es concluyente.

\subsection{Eje de tamaño: extrapolación a circuitos mayores}
\label{sec:res_tamano}

El entrenamiento llega hasta once qubits y la prueba hasta quince, así que los cuatro últimos
puntos de la Figura~\ref{fig:res_tamano} son extrapolación. Cabe destacar que Random Forest no
puede extrapolar por construcción.

% Figura ya exportada por 05c (v3_tamano_r2). La linea a trazos marca el ultimo tamaño que
% existe en entrenamiento.
\begin{figure}[htbp]
    \centering
    \caption{Coeficiente de determinación sobre el factor, por número de qubits. La línea a
             trazos marca el último tamaño presente en el entrenamiento.}
    \label{fig:res_tamano}

    \includegraphics[width=\textwidth]{cap08_v3/v3_tamano_r2}

    {\footnotesize Fuente: elaboración propia. Entre $69$ y $360$ circuitos por celda.}
\end{figure}

Dentro del rango conocido (5-11 qubits) los cinco modelos vuelven a moverse juntos, con Ridge algo por debajo y seperándose del resto en el observable correlador de vecinos. Lo interesante empieza a la derecha de la línea:

\begin{itemize}
    \item \textbf{La media de $X$ es donde la red de grafos gana.} A catorce y
          quince qubits saca a Random Forest $0{,}21$ y $0{,}23$ puntos de coeficiente. Es el único sitio de todo el trabajo donde la estructura del circuito sirve para predecir más que el resumen.
    \item \textbf{La media de $Z$ se hunde para todos.}
    \item \textbf{La paridad aguanta entera.} Se mantiene entre $0{,}89$ y $0{,}95$ hasta los
          quince qubits, y es el único observable que cruza la frontera sin sufrir, encabezado por Random Forest.
\end{itemize}

En resumen, \textbf{existe ventaja del grafo pero local}: aparece solo al extrapolar en tamaño, y solo en el observable donde queda señal que extrapolar.

\subsection{Eje de tipo: algoritmos nunca vistos}
\label{sec:res_tipo}

Es el eje más exigente de los tres. El entrenamiento es aleatorio al cien por cien, así que los
seis tipos con estructura del conjunto de prueba se evalúan a ciegas, y dentro de cada
tipo de observable.

% Figura ya exportada por 05c (v3_tipo_r2). Las celdas que se salen del eje van marcadas
% con un triangulo; el valor exacto esta en la tabla del cuaderno.
\begin{figure}[htbp]
    \centering
    \caption{Coeficiente de determinación sobre el factor, por tipo de circuito. Las celdas
             que se salen del eje van marcadas.}
    \label{fig:res_tipo}

    \includegraphics[width=\textwidth]{cap08_v3/v3_tipo_r2}

    {\footnotesize Fuente: elaboración propia. Entre $157$ y $1{.}115$ circuitos por celda.}
\end{figure}

Por lo general, \textbf{el coeficiente de determinación cae en la mayoría de los casos}. Y se destacan dos hallazgos

\begin{itemize}
    \item \textbf{En circuitos aleatorios (el caso conocido) todos funcionan}, entre
          $0{,}31$ y $0{,}57$, con Random Forest y las redes de grafos encabezan el primer puesto por
          centésimas.
    \item \textbf{La red de grafos gana en pocos casos} por enciam de Random Forest: predice mejor qft y bv en media X, en media de Y con random y tfim en correlación de vecinos.
    \item \textbf{En el resto de tipos con estructura nadie aprende gran cosa.} Random Forest gana por lo general en el resto. 
\end{itemize}

Sumando los tres ejes, el veredicto sobre la métrica que decide es \textbf{Random Forest}: gana
en validación, va en cabeza dentro del rango de tamaños conocido, y en el eje temporal gana con soltura. Las redes de grafos lo superan en dos sitios concretos:
grandes la media de $X$ al extrapolar en tamaño, y algunos circuitos especificos y de un observable concreto que no tienen relevancia. Entonces, se responde a la pregunta con la que se cerró el capítulo~\ref{cap:variable}:

\begin{center}
    \doublebox{
    \begin{minipage}{0.85\linewidth} \vspace{0.5em}
        \centering
        \textit{La estructura del circuito \textbf{sí} aporta información que el resumen
        agregado descarta, pero solo en la extrapolación del número de qubits. En el régimen conocido,
        treinta y cuatro columnas bastan.}
        \vspace{0.5em}
    \end{minipage}}
\end{center}


% =====================================================================================
\section{Resultados sobre la mitigación}
\label{sec:res_mitigacion}
% =====================================================================================

Este apartado mide la
distancia entre lo bien que prefice el factor y lo bien que corrige.
Por contrucción, un valor negativo significa que \textbf{habría sido mejor no tocar nada}.

% ⚠️ FORMATO: titulo arriba y fuente abajo (Instrucciones §1.3, p. 6).
% Cifras: 05d, tablas `mitigacion_mae` (columna «Sin mitigar») y `mitigacion_mejora`.
\begin{table}[htbp]
    \centering
    \caption{Mejora relativa del observable corregido frente a no corregir.}
    \label{tab:mitigacion}
    {\footnotesize
    Fuente: elaboración propia. \begin{tabular}{@{}lrrrrrrr@{}}
        \toprule
        & & \textbf{Sin mitigar} & \multicolumn{5}{c}{\textbf{Mejora relativa}} \\
        \cmidrule(l){4-8}
        \textbf{Observable} & \textbf{$n$} & \textbf{MAE} & \textbf{Ridge}
            & \textbf{RF} & \textbf{SAGE} & \textbf{GIN} & \textbf{GATv2} \\
        \midrule
        media $Z$    & $1{.}986$ & $0{,}0046$ & $-1{,}99$ & $\bm{+0{,}33}$ & $+0{,}14$ & $-0{,}02$ & $+0{,}06$ \\
        media $X$    & $3{.}393$ & $0{,}0177$ & $+0{,}50$ & $\bm{+0{,}81}$ & $+0{,}73$ & $+0{,}79$ & $+0{,}76$ \\
        media $Y$    & $1{.}096$ & $0{,}0044$ & $-1{,}49$ & $\bm{+0{,}37}$ & $+0{,}10$ & $+0{,}28$ & $+0{,}22$ \\
        paridad      & $1{.}825$ & $0{,}0111$ & $-1{,}91$ & $\bm{+0{,}73}$ & $+0{,}12$ & $+0{,}44$ & $-0{,}29$ \\
        corr. vec.   & $3{.}369$ & $0{,}0072$ & $-4{,}19$ & $+0{,}01$ & $-0{,}05$ & $\bm{+0{,}19}$ & $-0{,}46$ \\
        \bottomrule
    \end{tabular}}

    {\footnotesize Conjunto de prueba.}
\end{table}

% Figura ya exportada por 05d (mitigacion_mejora).
\begin{figure}[htbp]
    \centering
    \caption{Mejora relativa sobre no mitigar, por observable y modelo. El cero es
             exactamente no corregir.}
    \label{fig:res_mitigacion}

    \includegraphics[width=\textwidth]{cap08_mitigacion/mitigacion_mejora}

    {\footnotesize Fuente: elaboración propia. Conjunto de prueba.}
\end{figure}

El resultado que más dice del problema es el de Ridge. Su coeficiente sobre el factor era
\textbf{positivo en los cinco observables} (entre $0{,}41$ y $0{,}65$ en validación) y aun así,
aplicado, \textbf{empeora el resultado en cuatro de los cinco}. Por tanto, \textbf{un buen ajuste sobre el
factor no garantiza una buena mitigación}, y este trabajo lo mide en lugar de suponerlo.

En cuanto al resto. \textbf{Random Forest es el único que nunca empeora}, de $+0{,}01$ a
$+0{,}81$, y es el mejor en cuatro de los cinco observables. Entre las redes, GIN es la más
consistente y la única que aporta algo en el correlador de vecinos. Y la media de $X$ es el
único observable donde los cinco modelos ayudan, con mejoras de entre el $50$\,\% y el
$81$\,\%.

Queda un matiz que una media no puede dar. Medida como \textbf{proporción de circuitos que
mejoran}, Ridge corrige bien el $73$\,\% de las veces en la paridad y aun así su mejora media
es de $-1{,}91$. GIN mejora más circuitos ($93$\,\% frente a $89$\,\%) y sin
embargo su mejora media es menor ($+0{,}44$ frente a $+0{,}73$) de Random Forest, porque acierta más veces pero
falla más fuerte. Las dos métricas cuentan historias distintas.


Uniendo las dos mitades del capítulo, el trabajo termina con dos conclusiones que no coinciden.
Sobre la métrica que decide (la capacidad de predecir antes de ejecutar) el resumen tabular
basta en el régimen conocido y la estructura del grafo solo gana al extrapolar. Sobre la
aplicación, el orden es el mismo pero la distancia entre los modelos es mayor.


% =====================================================================================
% TABLA POR EJE — añadida sep-2026 a peticion del autor. Va DELANTE de la tabla resumen
% final a proposito: primero el desglose por eje, que es interpretable, y despues la
% mezcla de los 4.000, que no lo es tanto.
%
% 🔴 EJES AISLADOS, no ejes "a secas". Cada bloque fija los otros dos:
%      TIPO    -> estructurados y de HASTA 11 qubits  (1.816 circuitos)
%      TAMAÑO  -> aleatorios y de 12 a 15 qubits      (603)
%      TIEMPO  -> aleatorios y de hasta 11 qubits     (595)
%    Asi los tres son DISJUNTOS y una caida se puede atribuir a un eje. Sin aislar, los
%    2.802 estructurados y los 1.589 grandes comparten 986 circuitos y una caida en
%    «tipo» podria venir en realidad del tamaño. Es el mismo criterio con el que
%    §\ref{sec:res_tiempo} mide la deriva «en limpio».
%    ⚠️ Suman 3.014 de los 4.000: quedan fuera los 986 que son estructurados Y grandes.
%
% 🔴 LA MEDIA DE Y VA CON RAYA EN EL BLOQUE DE TIPO. No es un hueco ni un fallo: ningun
%    circuito estructurado define ese observable, asi que no hay ni una casilla que medir.
%
% PROCEDENCIA: evaluacion.evaluar() sobre logs/predicciones/pred_<modelo>_test.parquet,
% filtrando cada subconjunto. Mismo filtro que el resto del capitulo.
%
% ⚠️ NEGRITA = el mayor de la columna ANTES de redondear. En el bloque temporal casi todo
%    empata a dos decimales: eso NO es una errata, es el resultado.
% =====================================================================================
\begin{table}[htbp]
    \centering
    \caption{Coeficiente de determinación sobre el factor, por eje de generalización.}
    \label{tab:resumen_ejes}
    {\footnotesize
    Fuente: elaboración propia. \begin{tabular}{@{}lrrrrr@{}}
        \toprule
        & \textbf{media $Z$} & \textbf{media $X$} & \textbf{media $Y$}
        & \textbf{paridad} & \textbf{corr. vec.} \\
        \midrule
        \multicolumn{6}{@{}l}{\textit{Eje de tipo} — circuitos estructurados, hasta $11$ qubits} \\
        \midrule
        Ridge · $34$ variables & $-1{,}32$ & $-0{,}15$ & --- & $0{,}43$ & $-4{,}22$ \\
        Random Forest          & $\bm{-0{,}01}$ & $0{,}18$ & --- & $\bm{0{,}70}$ & $-0{,}09$ \\
        GNN · SAGE             & $-0{,}02$ & $0{,}09$ & --- & $-0{,}45$ & $-0{,}09$ \\
        GNN · GIN              & $-0{,}15$ & $\bm{0{,}24}$ & --- & $-0{,}23$ & $\bm{0{,}05}$ \\
        GNN · GATv2            & $-0{,}10$ & $-0{,}07$ & --- & $0{,}06$ & $-0{,}28$ \\
        \midrule
        \multicolumn{6}{@{}l}{\textit{Eje de tamaño} — circuitos aleatorios, de $12$ a $15$ qubits} \\
        \midrule
        Ridge · $34$ variables & $0{,}14$ & $0{,}35$ & $0{,}38$ & $0{,}60$ & $0{,}56$ \\
        Random Forest          & $\bm{0{,}30}$ & $0{,}37$ & $0{,}43$ & $\bm{0{,}82}$ & $0{,}51$ \\
        GNN · SAGE             & $0{,}17$ & $\bm{0{,}43}$ & $\bm{0{,}45}$ & $0{,}80$ & $0{,}62$ \\
        GNN · GIN              & $0{,}22$ & $0{,}36$ & $0{,}43$ & $0{,}74$ & $\bm{0{,}62}$ \\
        GNN · GATv2            & $0{,}19$ & $0{,}42$ & $0{,}42$ & $0{,}74$ & $0{,}61$ \\
        \midrule
        \multicolumn{6}{@{}l}{\textit{Eje temporal} — circuitos aleatorios, hasta $11$ qubits} \\
        \midrule
        Ridge · $34$ variables & $0{,}54$ & $0{,}54$ & $0{,}39$ & $0{,}68$ & $0{,}49$ \\
        Random Forest          & $0{,}53$ & $0{,}55$ & $0{,}39$ & $\bm{0{,}71}$ & $0{,}48$ \\
        GNN · SAGE             & $\bm{0{,}55}$ & $0{,}54$ & $0{,}39$ & $0{,}65$ & $\bm{0{,}49}$ \\
        GNN · GIN              & $0{,}54$ & $\bm{0{,}55}$ & $\bm{0{,}39}$ & $0{,}66$ & $0{,}49$ \\
        GNN · GATv2            & $0{,}54$ & $0{,}54$ & $0{,}38$ & $0{,}64$ & $0{,}49$ \\
        \bottomrule
    \end{tabular}}

    {\footnotesize Cada eje fija los otros dos, de modo que los
    tres subconjuntos son disjuntos y una caída se puede atribuir a un eje y no a la mezcla. La raya en la media de $Y$ del eje de
    tipo no es un hueco: ningún circuito estructurado define ese observable. En negrita, el
    mejor de cada columna antes de redondear.}
\end{table}

La Tabla~\ref{tab:resumen_final} cierra el capítulo con las dos métricas juntas, sobre el
conjunto de prueba entero y sin separar por ejes.

% =====================================================================================
% TABLA RESUMEN — añadida sep-2026 por indicación del tutor («una tabla al final a modo de
% resumen que compare todas las métricas y de todos los modelos»).
%
% ⚠️ FORMATO: titulo arriba y fuente abajo (Instrucciones §1.3, p. 6).
%
% 🔴 DOS BLOQUES Y NO UNA COLUMNA POR METRICA. Los cinco observables NO se promedian —viven
%    en escalas distintas—, asi que cada bloque conserva sus cinco columnas. Es la misma
%    regla que gobierna el resto del capitulo.
%
% 🔴 PROCEDENCIA DE LAS CIFRAS. Se calculan con `evaluacion.evaluar()` sobre los parquet de
%    logs/predicciones/pred_<modelo>_test.parquet, es decir, con EL MISMO filtro que el
%    resto del trabajo (casilla con señal y `f` real no negativa). Comprobado: el bloque de
%    mejora relativa reproduce exactamente la Tabla~\ref{tab:mitigacion}.
%    ⚠️ El R² sobre PRUEBA de GIN y GATv2 no aparecia antes en la memoria: los CSV de
%       logs/comparativa/comparativa_r2.csv solo traen hasta SAGE. Se calculo aqui.
%
% ⚠️ NEGRITA = el mayor de la columna ANTES de redondear. En la media de Y, SAGE (0,4231) y
%    Random Forest (0,4159) redondean los dos a 0,42 y solo uno va en negrita: no es una
%    errata. Dicho en el pie.
% =====================================================================================


\begin{table}[htbp]
    \centering
    \caption{Resumen de la comparativa: los ocho peldaños frente a las dos métricas, por
             observable, sobre el conjunto de prueba completo (sin distinción de eje).}
    \label{tab:resumen_final}
    {\footnotesize
    Fuente: elaboración propia. \begin{tabular}{@{}lrrrrr@{}}
        \toprule
        & \textbf{media $Z$} & \textbf{media $X$} & \textbf{media $Y$}
        & \textbf{paridad} & \textbf{corr. vec.} \\
        \midrule
        \multicolumn{6}{@{}l}{\textit{Coeficiente de determinación sobre el factor de
            supervivencia}} \\
        \midrule
        No mitigar             & $-0{,}17$ & $-0{,}23$ & $-0{,}30$ & $-0{,}82$ & $-0{,}20$ \\
        Factor constante       & $-0{,}48$ & $-1{,}02$ & $-0{,}16$ & $-0{,}40$ & $-1{,}26$ \\
        Ridge · $1$ variable   & $0{,}00$ & $0{,}09$ & $0{,}09$ & $0{,}72$ & $0{,}42$ \\
        Ridge · $34$ variables & $0{,}15$ & $0{,}41$ & $0{,}39$ & $0{,}84$ & $0{,}29$ \\
        Random Forest          & $\bm{0{,}42}$ & $0{,}45$ & $0{,}42$ & $\bm{0{,}90}$ & $0{,}55$ \\
        GNN · SAGE             & $0{,}36$ & $\bm{0{,}50}$ & $\bm{0{,}42}$ & $0{,}81$ & $0{,}61$ \\
        GNN · GIN              & $0{,}37$ & $0{,}48$ & $0{,}41$ & $0{,}81$ & $\bm{0{,}62}$ \\
        GNN · GATv2            & $0{,}36$ & $0{,}49$ & $0{,}40$ & $0{,}83$ & $0{,}60$ \\
        \midrule
        \multicolumn{6}{@{}l}{\textit{Mejora relativa sobre no mitigar}} \\
        \midrule
        No mitigar             & $0{,}00$ & $0{,}00$ & $0{,}00$ & $0{,}00$ & $0{,}00$ \\
        Factor constante       & $-4{,}77$ & $-1{,}00$ & $-2{,}75$ & $-0{,}76$ & $-5{,}63$ \\
        Ridge · $1$ variable   & $-1{,}82$ & $0{,}22$ & $-4{,}05$ & $-2{,}72$ & $-2{,}00$ \\
        Ridge · $34$ variables & $-1{,}99$ & $0{,}50$ & $-1{,}49$ & $-1{,}91$ & $-4{,}19$ \\
        Random Forest          & $\bm{0{,}33}$ & $\bm{0{,}81}$ & $\bm{0{,}37}$ & $\bm{0{,}73}$ & $0{,}01$ \\
        GNN · SAGE             & $0{,}14$ & $0{,}73$ & $0{,}10$ & $0{,}12$ & $-0{,}05$ \\
        GNN · GIN              & $-0{,}02$ & $0{,}79$ & $0{,}28$ & $0{,}44$ & $\bm{0{,}19}$ \\
        GNN · GATv2            & $0{,}06$ & $0{,}76$ & $0{,}22$ & $-0{,}29$ & $-0{,}46$ \\
        \bottomrule
    \end{tabular}}

    {\footnotesize Conjunto de prueba al completo ($4000$ circuitos), con el mismo filtro que
    el resto del capítulo. En negrita, el mejor de cada columna antes de redondear:
    en la media de $Y$ los dos primeros redondean a $0{,}42$ y solo uno está en negrita.}
\end{table}
